% Part: normal-modal-logic
% Chapter: tableaux
% Section: completeness

\documentclass[../../../include/open-logic-section]{subfiles}

\begin{document}

\olfileid{nml}{tab}{cpl}

\olsection{د \Log{K} بشپړتيا}

\begin{explain}
  د !!{tableau}s د طريقې د بشپړتيا د ښودلو لپاره بايد ثابته کړو
  چې که د $\Gamma \Proves !A$ د ښودلو لپاره تړلے !!{tableau}
  نۀ وي، نو $\Gamma \Entails/ !A$ وي، يعنې ضدمدل موجود وي۔
  خو «تړلے !!{tableau} نۀ شته» معنا دا ده چې د هغه د جوړولو
  هره ممکنه هڅه په تړلو کښې پاتې راځي۔ مهمه خبره دا ده چې که
  هره هڅه داسې پاتې راځي، نو يوه ټاکلې، \emph{منظمه او هر اړخيزه}
  طريقه هم څانګې نۀ شي تړلے۔ او که تړلے !!{tableau} موجود وي،
  همدغه منظمه او هر اړخيزه طريقه به يې وتړي۔ د دې يوې تابلو
  پرانيستې څانګې به د ضدمدل د تعريف لپاره ټول اړين معلومات ولري۔
  د داسې څانګې ضدمدل هغه ټول مخکښونه چې پر څانګه کارېدلي،
  د نړۍو په توګه اخلي؛ !!a{propositional
    variable}~$p$ په
  $\sigma$ کښې هله او يوازې هله رښتينے وي چې
  $\sFmla{\True}{p}[\sigma]$ په څانګه کښې راغلے وي۔
\end{explain}

\begin{defn}
  د !!a{tableau} يوه څانګه هغه مهال بشپړه بولو چې هرکله پکښې
  داسې مخکښ‌لرونکے !!{formula}~$\sFmla{S}{!A}[\sigma]$ وي چې
  قاعده پرې پلې کېدے شي، نو څانګه دا هم ولري:
  \begin{enumerate}
    \item د قضيېيزو يو-پر-بل اېښودونکو قاعدو په حالت کښې د
      قاعدې ټول اړوند پايلې، يعنې مخکښ‌لرونکي !!{formula}s؛
    \item د قضيېيزو څانګې وېشونکو قاعدو په حالت کښې له
      اړوندو پايلنو !!{formula}s څخه لږ تر لږه يو؛
    \item د هغو مودالي قاعدو په حالت کښې چې نوے مخکښ غواړي،
      لږ تر لږه يوه ممکنه پايله؛
    \item د هغو مودالي قاعدو په حالت کښې چې کارېدلے مخکښ
      غواړي، د څانګې د هر پلي کېدونکي کارېدلي مخکښ لپاره
      اړونده پايله۔
  \end{enumerate}
  په منجمده سرچينه کښې وروستۍ ماده «هر راغلے مخکښ» وايي؛
  د قاعدې له مخې يوازې د مخکښ په غځېدلې بڼه پلي کېدونکي
  کارېدلي مخکښونه مراد دي۔
\end{defn}

\begin{explain}
د بېلګې په توګه، که بشپړه څانګه $\sFmla{\True}{!B
  \land !C}[\sigma]$ ولري، نو $\sFmla{\True}{!B}[\sigma]$
او $\sFmla{\True}{!C}[\sigma]$ هم لري۔ که
$\sFmla{\True}{!B \lor !C}[\sigma]$ ولري، نو لږ تر لږه
$\sFmla{\True}{!B}[\sigma]$ يا $\sFmla{\True}{!C}[\sigma]$
لري۔ په منجمده سرچينه کښې د لومړۍ مقدمې مخکښ پاتې شوے
او د دويمې بېلګې لومړۍ پايله په تېروتنه د B د ناسموالي
نښه ده؛ دلته د قضيېيزو قاعدو رښتينې پايلې ښودل شوې دي۔
که څانګه \iftag{prvBox}
{$\sFmla{\False}{\Box !A}[\sigma]$} {$\sFmla{\True}{\Diamond !A}[\sigma]$}
ولري، نو د لږ تر لږه يوه~$n$ لپاره
\iftag{prvBox} {$\sFmla{\False}{!A}[\sigma.n]$}
{$\sFmla{\True}{!A}[\sigma.n]$} هم لري۔ او که څانګه
\iftag{prvBox} {$\sFmla{\True}{\Box !A}[\sigma]$}
{$\sFmla{\False}{\Diamond !A}[\sigma]$} ولري، نو د هر هغه~$n$
لپاره چې $\sigma.n$ په څانګه کښې کارېدلے وي،
\iftag{prvBox} {$\sFmla{\True}{!A}[\sigma.n]$}
{$\sFmla{\False}{!A}[\sigma.n]$} هم لري۔ په منجمده سرچينه
کښې د دې څلورو مقدماتو او څلورو پايلو دنننے فارمول
غورځېدلے؛ دلته د مخکنيو مودالي قاعدو سره سم راوړل شوے دے۔
\end{explain}

\begin{prop}\ollabel{prop:complete-tableau}
  هر متناهي~$\Gamma$ داسې !!a{tableau} لري چې هره څانګه ئې بشپړه وي۔
\end{prop}

\begin{proof}
  د~$\Gamma$ د !!a{tableau} يوه پرانيستې څانګه واخلئ۔ پر دې
  څانګه يوازې متناهي ډېر داسې مخکښ‌لرونکي !!{formula}s دي چې
  قاعده پرې پلې کېدے شي۔ په يوه ټاکلي ترتيب (مثلاً له پاسه
  تر ښکته) د دغو !!{formula}s له هر يوه لپاره چې پورته (۱)–(۴) شرطونه
  يې لا نۀ وي پوره کړي، پلې کېدونکې قاعدې وکاروئ او څانګه
  وغځوئ۔ کله کله څانګه وېشل کېږي؛ د هرې نوې څانګې په سر کښې
  پر پاتې مخکښ‌لرونکو !!{formula}s همدغه قاعدې پلې کړئ۔
  ځکه چې پر څانګه د مخکښ‌لرونکو !!{formula}s او کارېدلو
  مخکښونو شمېر متناهي دے، دا کړنلاره په پايله کښې اصلي
  څانګه په (کېدے شي ګڼو) څانګو غځوي۔ په هره يوه يې همدا
  کړنلاره بيا پلې کړئ او تکرار يې کړئ۔ د جوړښت له مخې به
  هره څانګه بشپړه وي۔ په منجمده سرچينه کښې دلته «تړلې»
  ليکل شوي، خو دا د همدې قضیې له غوښتنې سره نۀ جوړېږي۔
\end{proof}

\begin{thm}[بشپړتيا]
  \ollabel{thm:tableau-completeness}
  که~$\Gamma$ تړلے !!{tableau} ونۀ لري، نو~$\Gamma$ د صدق وړ دے۔
\end{thm}

\begin{proof}
د پورته قضيې له مخې~$\Gamma$ داسې !!a{tableau} لري چې هره
څانګه ئې بشپړه ده۔ ځکه تړلے !!{tableau} نۀ لري، نو په دغه
!!a{tableau} کښې لږ تر لږه يوه پرانيستې او بشپړه څانګه شته۔
$\Delta$ دې پر څانګه د مخکښ‌لرونکو !!{formula}s سټ وي، او
$P(\Delta)$ دې پکښې د راغلو مخکښونو سټ وي۔

مدل~$\mModel{M(\Delta)} = \tuple{P(\Delta), R, V}$ داسې تعريفوو
چې نړۍې ئې په~$\Delta$ کښې راغلي مخکښونه دي، او د لاسرسي
اړيکه يې داسې ده:
\[
R\sigma\sigma' \quad \text{هله او يوازې هله} \quad
\sigma'=\sigma.n \quad \text{د کوم~$n$ لپاره}
\]
او
\[
V(p) = \Setabs{\sigma}{\sFmla{\True}{p}[\sigma] \in \Delta}.
\]
پر~$!A$ د استقرا له لارې ښيو چې که
$\sFmla{\True}{!A}[\sigma] \in \Delta$ وي، نو
$\mSat{M(\Delta)}{!A}[\sigma]$ وي؛ او که
$\sFmla{\False}{!A}[\sigma] \in \Delta$ وي، نو
$\mSat/{M(\Delta)}{!A}[\sigma]$ وي۔

\begin{enumerate}
  \item \indcase{!A}{p}{که
    $\sFmla{\True}{\indfrm}[\sigma] \in \Delta$ وي، نو د~$V$
    د تعريف له مخې $\sigma \in V(p)$ وي؛ له دې امله
    $\mSat{M(\Delta)}{\indfrm}[\sigma]$ وي۔

    که $\sFmla{\False}{\indfrm}[\sigma] \in \Delta$ وي، نو
    $\sFmla{\True}{\indfrm}[\sigma] \notin \Delta$ وي؛ ځکه
    که داسې نۀ وي، څانګه به تړلې وي۔ نو $\sigma \notin V(p)$
    او له دې امله $\mSat/{M(\Delta)}{\indfrm}[\sigma]$ وي۔}
  \item \indcase{!A}{\lnot !B}{\iftag{probNot}{تمرين۔}{که
      $\sFmla{\True}{\indfrm}[\sigma] \in \Delta$ وي، نو د
      څانګې د بشپړتيا له مخې
      $\sFmla{\False}{!B}[\sigma] \in \Delta$ وي۔ د استقرا د
      فرضيې له مخې $\mSat/{M(\Delta)}{!B}[\sigma]$ وي، نو
      $\mSat{M(\Delta)}{\indfrm}[\sigma]$ وي۔

      که $\sFmla{\False}{\indfrm}[\sigma] \in \Delta$ وي، نو د
      څانګې د بشپړتيا له مخې
      $\sFmla{\True}{!B}[\sigma] \in \Delta$ وي۔ د استقرا د
      فرضيې له مخې $\mSat{M(\Delta)}{!B}[\sigma]$ وي، نو
      $\mSat/{M(\Delta)}{\indfrm}[\sigma]$ وي۔}}
  \item \indcase{!A}{!B \land !C}{\iftag{probAnd}{تمرين۔}{که
      $\sFmla{\True}{\indfrm}[\sigma] \in \Delta$ وي، نو د
      څانګې د بشپړتيا له مخې دواړه
      $\sFmla{\True}{!B}[\sigma] \in \Delta$ او
      $\sFmla{\True}{!C}[\sigma] \in \Delta$ وي۔ د استقرا د
      فرضيې له مخې $\mSat{M(\Delta)}{!B}[\sigma]$ او
      $\mSat{M(\Delta)}{!C}[\sigma]$ وي؛ نو
      $\mSat{M(\Delta)}{\indfrm}[\sigma]$ وي۔

      که $\sFmla{\False}{\indfrm}[\sigma] \in \Delta$ وي، نو د
      څانګې د بشپړتيا له مخې يا
      $\sFmla{\False}{!B}[\sigma] \in \Delta$ يا
      $\sFmla{\False}{!C}[\sigma] \in \Delta$ وي۔ د استقرا د
      فرضيې له مخې يا $\mSat/{M(\Delta)}{!B}[\sigma]$ يا
      $\mSat/{M(\Delta)}{!C}[\sigma]$ وي؛ نو
      $\mSat/{M(\Delta)}{\indfrm}[\sigma]$ وي۔ په منجمده
      سرچينه کښې د دويم غړي پر ځاے لومړنے بيا ليکل شوے و۔}}
  \item \indcase{!A}{!B \lor !C}{\iftag{probOr}{تمرين۔}{که
      $\sFmla{\True}{\indfrm}[\sigma] \in \Delta$ وي، نو د
      څانګې د بشپړتيا له مخې يا
      $\sFmla{\True}{!B}[\sigma] \in \Delta$ يا
      $\sFmla{\True}{!C}[\sigma] \in \Delta$ وي۔ د استقرا د
      فرضيې له مخې يا $\mSat{M(\Delta)}{!B}[\sigma]$ يا
      $\mSat{M(\Delta)}{!C}[\sigma]$ وي؛ نو
      $\mSat{M(\Delta)}{\indfrm}[\sigma]$ وي۔

      که $\sFmla{\False}{\indfrm}[\sigma] \in \Delta$ وي، نو د
      څانګې د بشپړتيا له مخې دواړه
      $\sFmla{\False}{!B}[\sigma] \in \Delta$ او
      $\sFmla{\False}{!C}[\sigma] \in \Delta$ وي۔ د استقرا د
      فرضيې له مخې $\mSat/{M(\Delta)}{!B}[\sigma]$ او
      $\mSat/{M(\Delta)}{!C}[\sigma]$ وي؛ نو
      $\mSat/{M(\Delta)}{\indfrm}[\sigma]$ وي۔ په منجمده
      سرچينه کښې د دويم غړي پر ځاے لومړنے بيا ليکل شوے و۔}}
  \item \indcase{!A}{!B \lif !C}{\iftag{probIf}{تمرين۔}{که
      $\sFmla{\True}{\indfrm}[\sigma] \in \Delta$ وي، نو د
      څانګې د بشپړتيا له مخې يا
      $\sFmla{\False}{!B}[\sigma] \in \Delta$ يا
      $\sFmla{\True}{!C}[\sigma] \in \Delta$ وي۔ د استقرا د
      فرضيې له مخې يا $\mSat/{M(\Delta)}{!B}[\sigma]$ يا
      $\mSat{M(\Delta)}{!C}[\sigma]$ وي؛ نو
      $\mSat{M(\Delta)}{\indfrm}[\sigma]$ وي۔

      که $\sFmla{\False}{\indfrm}[\sigma] \in \Delta$ وي، نو د
      څانګې د بشپړتيا له مخې دواړه
      $\sFmla{\True}{!B}[\sigma] \in \Delta$ او
      $\sFmla{\False}{!C}[\sigma] \in \Delta$ وي۔ د استقرا د
      فرضيې له مخې $\mSat{M(\Delta)}{!B}[\sigma]$ او
      $\mSat/{M(\Delta)}{!C}[\sigma]$ وي؛ نو
      $\mSat/{M(\Delta)}{\indfrm}[\sigma]$ وي۔ په منجمده
      سرچينه کښې د دويم غړي پر ځاے لومړنے بيا ليکل شوے و۔}}
  \item \indcase{!A}{\Box !B}{\iftag{probBox}{تمرين۔}{که
      $\sFmla{\True}{\indfrm}[\sigma] \in \Delta$ وي، نو د
      څانګې د بشپړتيا له مخې د هر کارېدلي~$\sigma.n$ لپاره
      $\sFmla{\True}{!B}[\sigma.n] \in \Delta$ وي؛ يعنې د هر
      هغه~$\sigma' \in P(\Delta)$ لپاره چې $R\sigma\sigma'$ وي۔
      د استقرا د فرضيې له مخې د هر هغه~$\sigma'$ لپاره چې
      $R\sigma\sigma'$ وي،
      $\mSat{M(\Delta)}{!B}[\sigma']$ وي؛ نو
      $\mSat{M(\Delta)}{\indfrm}[\sigma]$ وي۔

      که $\sFmla{\False}{\indfrm}[\sigma] \in \Delta$ وي، نو د
      څانګې د بشپړتيا له مخې د کوم~$\sigma.n$ لپاره
      $\sFmla{\False}{!B}[\sigma.n] \in \Delta$ وي۔ د استقرا د
      فرضيې له مخې $\mSat/{M(\Delta)}{!B}[\sigma.n]$ وي۔
      ځکه چې $R\sigma(\sigma.n)$، داسې~$\sigma'$ شته چې
      $\mSat/{M(\Delta)}{!B}[\sigma']$ وي؛ نو
      $\mSat/{M(\Delta)}{\indfrm}[\sigma]$ وي۔}}
  \item \indcase{!A}{\Diamond !B}{\iftag{probDiamond}{تمرين۔}{که
      $\sFmla{\True}{\indfrm}[\sigma] \in \Delta$ وي، نو د
      څانګې د بشپړتيا له مخې د کوم~$\sigma.n$ لپاره
      $\sFmla{\True}{!B}[\sigma.n] \in \Delta$ وي۔ د استقرا د
      فرضيې له مخې $\mSat{M(\Delta)}{!B}[\sigma.n]$ وي۔
      ځکه چې $R\sigma(\sigma.n)$، داسې~$\sigma'$ شته چې
      $\mSat{M(\Delta)}{!B}[\sigma']$ وي؛ نو
      $\mSat{M(\Delta)}{\indfrm}[\sigma]$ وي۔

      که $\sFmla{\False}{\indfrm}[\sigma] \in \Delta$ وي، نو د
      څانګې د بشپړتيا له مخې د هر کارېدلي~$\sigma.n$ لپاره
      $\sFmla{\False}{!B}[\sigma.n] \in \Delta$ وي؛ يعنې د هر
      هغه~$\sigma' \in P(\Delta)$ لپاره چې $R\sigma\sigma'$ وي۔
      د استقرا د فرضيې له مخې د هر هغه~$\sigma'$ لپاره چې
      $R\sigma\sigma'$ وي،
      $\mSat/{M(\Delta)}{!B}[\sigma']$ وي؛ نو
      $\mSat/{M(\Delta)}{\indfrm}[\sigma]$ وي۔}}
\end{enumerate}
ځکه چې $\Gamma \subseteq \Delta$، نو
$\mSat{M(\Delta)}{\Gamma}$ وي۔
\end{proof}

\begin{prob}
د \olref[nml][tab][cpl]{thm:tableau-completeness} ثبوت بشپړ کړئ۔
\end{prob}

\begin{cor}
\ollabel{cor:entailment-completeness}
که $\Gamma \Entails !A$ وي، نو $\Gamma \Proves !A$ وي۔
\end{cor}

\begin{cor}
\ollabel{cor:weak-completeness}
که~$!A$ په ټولو مدلونو کښې رښتينے وي، نو $\Proves !A$ وي۔
\end{cor}

\end{document}
